\documentclass[10pt,letterpaper]{article}
\usepackage[margin=0.8in]{geometry}
\usepackage[T1]{fontenc}
\usepackage{lmodern}
\usepackage{amsmath,amssymb,booktabs,graphicx,float}
\usepackage[hidelinks]{hyperref}
\hypersetup{pdftitle={COMP 766 assignment 1: Curve inference with relaxation labeling}}
\graphicspath{{assets/}}
\setlength{\parindent}{0pt}
\setlength{\parskip}{5pt}
\newcommand{\asset}[2]{\includegraphics[width=#1\linewidth]{#2}}
\title{COMP 766 assignment 1\\Curve inference with relaxation labeling}
\author{} % Add your name and student ID here.
\date{October 9, 2026}
\begin{document}
\maketitle
\vspace{-2em}

\section{Approach}
This report applies the curve inference method of Parent and
Zucker~\cite{parent1989} using the relaxation framework of Hummel and
Zucker~\cite{hummel1983}. Each image is converted to grayscale and convolved
with a Gaussian-based line template at 16 rotations. The responses give
initial tangent confidences. Cocircularity and curvature consistency measure
how well neighboring tangents support each other, and radial relaxation
updates the confidences.

The three images are spaghetti, fingerprint, and hair. Spaghetti has clear
strands and crossings. Fingerprint has fine, low-contrast ridges that change
direction. Hair combines strands of different widths with highlights,
shadows, and closely spaced parallel curves.

\begin{table}[!htbp]
\centering\small
\begin{tabular}{ll}
\toprule
Parameter & Value \\
\midrule
Kernel rotations & $\theta_n=n\cdot \pi/16$, $n=0,\ldots,15$ \\
Kernel sampling & $31\times31$ on $[-2,2]\times[-2.5,2.5]$ \\
Gaussian widths & $\sigma_y=1$, $\sigma_1=0.8$, $\sigma_2=0.3$, $\sigma_3=0.8$ \\
Gaussian weights & $A=0.5$, $B=1.5$, $C=0.5$ \\
Initialization & Softmax temperature $\tau=0.005$ \\
Support & Radius $r=5$, seven curvature classes, $c_{\min}=0.1$ \\
Curvature range & $[-0.2,0.2]$ inverse pixels \\
Relaxation & Five updates, step $\eta=1$, tolerance $10^{-4}$ \\
Support reference & $s_{\min}=3$, $s_{\max}=10$ \\
Overlay thresholds & Initial $p>0.20$ and $p>0.50$, final $p>0.50$ \\
Overlay spacing & 4 pixels in full views, 2 pixels in crops \\
\bottomrule
\end{tabular}
\caption{Parameters used for all three images.}
\label{tab:parameters}
\end{table}

\section{Part I: initial tangent estimates}
\subsection{Measurement model and normalization}
The line template combines three Gaussian terms across the line, weighted by
$A$, $-B$, and $C$, with a Gaussian envelope along it. We sample the template
on a $31\times31$ grid. For each $\theta_n$, the sampling coordinates are rotated as
$(x',y')=(x\cos\theta_n+y\sin\theta_n,-x\sin\theta_n+y\cos\theta_n)$.
Each kernel is divided by the sum of its sampled values, giving it a total
weight of one. Convolution uses symmetric padding, which mirrors pixels at
the image edges instead of introducing a black border.

If $f_i(n)$ denotes the response for orientation $n$, initialization is
\begin{equation}
p_i^{(0)}(n)=\frac{\exp(f_i(n)/\tau)}{\sum_{m=0}^{15}\exp(f_i(m)/\tau)}.
\label{eq:softmax}
\end{equation}
The temperature $\tau=0.005$ controls how strongly the largest responses
dominate the initial confidences. The code uses SciPy's softmax function.

\subsection{Orientation in image coordinates}
The reference kernel is vertical. Its tangent in continuous sampling coordinates
is proportional to $(-\sin\theta,\cos\theta)$. Horizontal and vertical sample
spacings are $\Delta x=4/30$ and $\Delta y=5/30$, so its image-pixel angle is
\begin{equation}
\phi_n=\operatorname{atan2}\!\left(\frac{\cos\theta_n}{\Delta y},
                              \frac{-\sin\theta_n}{\Delta x}\right)\bmod\pi.
\label{eq:angles}
\end{equation}
Image rows increase downward. Both the displayed vectors and the support
calculation use $\phi_n$. Because the horizontal and vertical sample spacings
differ, equally spaced kernel rotations produce unevenly spaced tangent
angles in image pixels.

\begin{figure}[!htbp]
\centering
\asset{0.32}{spaghetti_initial.png}\hfill
\asset{0.32}{fingerprint_initial.png}\hfill
\asset{0.32}{hair_initial.png}
\caption{Initial estimates for spaghetti, fingerprint, and hair, from left to right. Blue vectors mark $p>0.20$ at spacing 4 pixels. The enlarged views in Figure~\ref{fig:details} make competing directions easier to inspect.}
\label{fig:initial}
\end{figure}

\begin{figure}[!htbp]
\centering
\asset{0.32}{spaghetti_initial_p05.png}\hfill
\asset{0.32}{fingerprint_initial_p05.png}\hfill
\asset{0.32}{hair_initial_p05.png}
\caption{Initial estimates at $p>0.50$ for spaghetti, fingerprint, and hair, from left to right. The probabilities and 4-pixel spacing are the same as in Figure~\ref{fig:initial}.}
\label{fig:initial-p05}
\end{figure}

\subsection{Initial results}
Increasing the display threshold from $p>0.20$ (Figure~\ref{fig:initial}) to
$p>0.50$ (Figure~\ref{fig:initial-p05}) reduces the displayed vectors from
8,685 to 2,794 for spaghetti, 14,319 to 790 for fingerprint, and 17,979 to
2,689 for hair, removing many competing directions.
Spaghetti retains vectors along most strands. Fingerprint and hair become
sparser, with gaps where no orientation exceeds the higher threshold.
Since the initial probabilities sum to one at each pixel, $p>0.50$ permits at
most one displayed orientation per sampled pixel. Both thresholds are shown
to distinguish the competing initial directions from the strongest responses.

Spaghetti's initial vectors mostly follow the strands, but crossings and
endpoints have many conflicting directions. The template's negative center
also produces responses along boundaries and darker parts of the bright
strands. Fingerprint and hair contain more conflicting estimates. In the
fingerprint crop, some vectors follow the vertical ridges while others point
across them. Many hair vectors follow the main diagonal direction, with
additional vectors pointing across the strands.

\section{Part II: cocircularity and curvature consistency}
\subsection{Cocircularity compatibility}
Two tangents are cocircular if the same circle is tangent to both at their
positions~\cite{parent1989}. Tangents on the same straight line are the
limiting case of infinite radius. Here $\lambda$ and $\lambda'$ label the
tangent directions. To match the paper's notation, their angles are written
as $\theta_\lambda$ and $\theta_{\lambda'}$, using the image angles
$\phi_\lambda$ and $\phi_{\lambda'}$ from Part I.

For pixels $(x_i,y_i)$ and $(x_j,y_j)$, define
\begin{equation}
\Delta x=x_j-x_i,\qquad \Delta y=y_j-y_i,\qquad
d_{ij}=\sqrt{\Delta x^2+\Delta y^2},\qquad
\theta_{ij}=\arctan(\Delta y/\Delta x).
\end{equation}
The angle $\theta_{ij}$ describes the line joining the pixels, as in
the paper's Equation (5.1). The code evaluates it using
$\operatorname{atan2}(\Delta y,\Delta x)\bmod\pi$ to handle vertical lines.
The paper's signed interior-angle function is
\begin{equation}
\Gamma(\beta,\gamma)=
\begin{cases}
\gamma-\beta, & |\gamma-\beta|\leq\pi/2,\\
\gamma-\beta-\pi, & \pi/2<\gamma-\beta\leq\pi,\\
\gamma-\beta+\pi, & -\pi\leq\gamma-\beta<-\pi/2.
\end{cases}
\end{equation}
Exact cocircularity requires equal interior angles,
$\Gamma(\theta_\lambda,\theta_{ij})=
\Gamma(\theta_{ij},\theta_{\lambda'})$.
For quantized positions and orientations, the implementation uses the paper's
discrete condition, Equation (5.5):
\begin{equation}
\left|\Gamma(\theta_\lambda,\theta_{ij})-
\Gamma(\theta_{ij},\theta_{\lambda'})\right|
<\epsilon+2\alpha,\qquad
\epsilon=\frac{\pi}{16},\quad \alpha=\arcsin(1/d_{ij}).
\label{eq:cocircularity}
\end{equation}
The tolerance $\epsilon$ accounts for sampled orientations. The term
$2\alpha$ accounts for uncertainty in the joining line when each tangent's
position can vary within half a pixel. The target pixel is excluded from its
own neighborhood, so $d_{ij}\geq1$.

The paper's Figure 9 and Equation (5.6) reduce compatibility gradually
outside the accepted angular range. The implementation simplifies this to
a step function:
\begin{equation}
c_{ij}(\lambda,\lambda')=
\begin{cases}
1, & \text{if Equation~\eqref{eq:cocircularity} holds},\\
c_{\min}=0.1, & \text{otherwise}.
\end{cases}
\end{equation}
Every pair outside the tolerance receives compatibility $0.1$, regardless
of its angular error. This gives those pairs weak support without the
paper's gradual transition.

\subsection{Curvature classes and support}
Curvature classes separate support for nearly straight curves from support for
curves bending in either direction. For a tangent with orientation $\phi$ at
pixel $i$, consider the circle tangent to it that passes through neighbor $j$.
The implementation computes its signed curvature as
\begin{equation}
\kappa_{ij}(\phi)=
\frac{2(\cos\phi\,\Delta y-\sin\phi\,\Delta x)}{d_{ij}^2}.
\label{eq:curvature}
\end{equation}
Its magnitude is the reciprocal of the circle's radius; zero curvature means
a straight line. This formula follows from the radius estimate in the paper's
Equation (5.7), with a sign assigned according to the tangent and displacement
directions used in the code.

The implementation uses seven curvature classes with the following boundaries,
computed from its neighborhood radius and maximum-curvature defaults:
\begin{equation}
(-0.20,-0.16,-0.12,-0.08,0.08,0.12,0.16,0.20).
\end{equation}
These eight boundaries define one central class and three classes on each side.
The central class, $[-0.08,0.08)$, groups straight lines and gentle bends.
Its width comes from the heuristic
$\min(2/r^2,\kappa_{\max}/2)$, with $r=5$ and $\kappa_{\max}=0.2$.
The term $2/r^2$ is the curvature given by Equation~\eqref{eq:curvature}
when a neighbor at distance $r$ lies one pixel away from the tangent line.
The other classes divide each remaining range into equal intervals of width
$0.04$. Curvatures outside $[-0.2,0.2]$ are excluded, so circles with radii
below 5 pixels provide no support. Each interval includes its lower boundary
and excludes its upper boundary, except that $0.2$ is included in the last class.

For each tangent, support is summed separately within each class.
A neighboring tangent contributes only if its previously selected class also
contains the target pixel. This is the curvature-consistency check in
Section V-D: both tangents must agree that the other lies in their expected
curve region.

Let $b_{ij}(\lambda)$ denote the class implied by tangent $\lambda$ at $i$
and neighbor $j$, and let $k_j(\lambda')$ denote the previously selected class
of tangent $\lambda'$ at $j$. The support is
\begin{align}
S_{ik}(\lambda)&=\sum_{j\in N_r(i)}\sum_{\lambda'}
 c_{ij}(\lambda,\lambda')\,
 \mathbf1\{b_{ij}(\lambda)=k\}\,
 \mathbf1\{b_{ji}(\lambda')=k_j(\lambda')\geq0\}\,p_j(\lambda'),\\
s_i(\lambda)&=\max_k S_{ik}(\lambda),\qquad
k_i(\lambda)=\operatorname*{arg\,max}_k S_{ik}(\lambda).
\label{eq:support}
\end{align}
The first indicator places the neighbor in class $k$ of the target tangent.
The second checks that the target belongs to the neighbor's selected class.
The largest class support becomes $s_i(\lambda)$, and its class becomes
$k_i(\lambda)$. For the reverse check, the code swaps $i$ and $j$, reversing
the displacement and hence the curvature sign. It therefore checks membership
in both directions rather than requiring identical class indices.

The first pass selects classes without the curvature-consistency check because
no previous estimates exist. A second pass uses those estimates to compute
consistent support. During relaxation, each pass uses the preceding pass's
classes. Excluded curvatures and tangents with zero support receive class $-1$.

\subsection{Local support neighborhood}
Each pixel receives support from the 80 other pixel positions within a
radius of 5 pixels. Neighbors outside the image are skipped, leaving fewer
supporting neighbors at image boundaries.

The code precomputes compatibility coefficients and curvature classes for
each neighbor offset. Matrix operations sum the neighboring confidences,
and processing blocks of rows limits memory use.

\subsection{Support results}
\begin{figure}[!htbp]
\centering
\asset{0.64}{spaghetti_support.png}\\[3pt]
\asset{0.64}{fingerprint_support.png}\\[3pt]
\asset{0.64}{hair_support.png}
\caption{Spaghetti, fingerprint, and hair, from top to bottom. Left panels show supported initial candidates; right panels show maximum raw support across orientations. Display thresholds are 0.18, 0.44, and 0.61 of each image's global maximum support.}
\label{fig:support}
\end{figure}

For the support overlays, each image's support values are divided by their
global maximum. Only candidates with initial confidence $p>0.20$ are eligible
for display. The display threshold is the 90th percentile of the normalized
maximum support at each pixel, rounded to two decimal places. This scaling
is used for visualization; relaxation uses raw support.

At these thresholds, 6,231 of 8,685 sampled spaghetti candidates remain,
or 71.7\%. Most vectors along the strands remain, while many vectors pointing
across them at intersections disappear. Fingerprint retains 2,213 of 14,319
candidates, or 15.5\%. Its central crop loses many vectors across the ridges,
but also loses some plausible ridge directions. Hair retains 3,243 of 17,979
candidates, or 18.0\%, mostly along the main strand direction.

The retained counts depend on both support and the display thresholds.
The overlays also retain vectors along the finger boundary, shadows, and
neighboring parallel strands, where local agreement can occur outside the
curves of interest.

\section{Part III: radial relaxation labeling}
\subsection{Update rule and reference scale}
The support score is
\begin{equation}
A(p)=\sum_i\sum_\lambda p_i(\lambda)s_i(\lambda).
\end{equation}
The plots show $A(p)/N$, where $N$ is the number of image pixels, to account
for differences in image size. Before each update, raw support is converted
to signed evidence:
\begin{equation}
q_i(\lambda)=\frac{s_i(\lambda)-s_{\min}}{s_{\max}-s_{\min}}.
\end{equation}
The reference $s_{\max}=2r=10$ comes from a horizontal line with confidence
one, which has ten supporting neighbors in this neighborhood. Following the
endpoint argument in Equation (6.15) of Parent and Zucker~\cite{parent1989},
a minimum line confidence of $0.6$ gives $s_{\min}=0.6s_{\max}/2=3$.
Raw support can exceed $s_{\max}$, and the code does not clip $q$ to $[-1,1]$.

Each orientation is updated independently as a pair: tangent confidence $p$
and no-line confidence $1-p$. Their evidence values are $q$ and $-q$.
Subtracting the smaller evidence value from both gives
$[2\max(q,0),2\max(-q,0)]$. Adding $\eta$ times these values to the confidence
pair and dividing by its new sum gives the radial rule from Appendix A of
the trace inference paper:
\begin{equation}
p^{(t+1)}=\frac{p^{(t)}+2\eta\max(q,0)}{1+2\eta|q|}.
\label{eq:update}
\end{equation}
Positive $q$ raises the tangent confidence, negative $q$ lowers it, and
$q=0$ leaves it unchanged. Each confidence stays in $[0,1]$. After
initialization, orientations no longer share a sum-to-one constraint, so
multiple directions can survive at a crossing.

After each update, the code recomputes support using the preceding curvature
classes and selects new classes. The run stops after five updates or when
the largest absolute confidence change is at most $10^{-4}$. All three
images reached the five-update limit.

\begin{figure}[!htbp]
\centering
\asset{1.0}{spaghetti_relaxation.png}
\caption{Spaghetti tangent estimates before relaxation (left, $p>0.20$) and after five updates (right, $p>0.50$).}
\label{fig:relaxation}
\end{figure}

\begin{figure}[!htbp]
\centering
\asset{1.0}{fingerprint_relaxation.png}
\caption{Fingerprint tangent estimates before relaxation (left, $p>0.20$) and after five updates (right, $p>0.50$).}
\label{fig:fingerprint-relaxation}
\end{figure}

\begin{figure}[!htbp]
\centering
\asset{1.0}{hair_relaxation.png}
\caption{Hair tangent estimates before relaxation (left, $p>0.20$) and after five updates (right, $p>0.50$).}
\label{fig:hair-relaxation}
\end{figure}

\begin{figure}[!htbp]
\centering\asset{0.65}{support_history.png}
\caption{Average local support per pixel over five updates for spaghetti, fingerprint, and hair.}
\label{fig:history}
\end{figure}

\begin{minipage}{\linewidth}
\section{Results after five updates}
\subsection{Support histories}
Spaghetti's support per pixel falls from 1.049 to 0.820 on the first update, then
rises to 2.270 by update five. Fingerprint falls slightly from 1.391 to 1.363,
then reaches 5.192. Hair increases from 1.610 to 7.005. The final scores are
approximately 2.16, 3.73, and 4.35 times their initial values, respectively.
\end{minipage}

Five updates were used as a practical limit. Support is still rising in
Figure~\ref{fig:history}, and the largest confidence changes on the fifth
update are approximately $0.51$, $0.51$, and $0.43$ for spaghetti, fingerprint,
and hair. None of the runs meets the $10^{-4}$ stopping tolerance.

\subsection{Local behavior}
\begin{figure}[!htbp]
\centering
\asset{0.78}{spaghetti_detail.png}\\[3pt]
\asset{0.78}{fingerprint_detail.png}\\[3pt]
\asset{0.78}{hair_detail.png}
\caption{Initial, supported, and relaxed central crops. Rows are spaghetti, fingerprint, and hair. The $128\times128$ regions start at $(536,348)$, $(192,192)$, and $(256,256)$ in column-row coordinates, respectively. Vectors are sampled every 2 pixels.}
\label{fig:details}
\end{figure}

Spaghetti shows the clearest visual improvement. Vectors follow the strands
more consistently, and many conflicting directions disappear. Both strand
directions remain visible near the central crossing. Several parallel rows
of vectors still span each strand's width because the method does not select
a single centerline or determine which strand passes over the other.

Fingerprint's vectors follow the upper and right-hand ridges more consistently,
but the central crop still has gaps and isolated horizontal fragments. Hair's
vectors increasingly follow the main diagonal direction, yet spread across
wide bands and some dark gaps between strands.

\section{Limitations of Assumptions}

Recovering a bifurcation would require linking an incoming curve to its
outgoing branches. The implementation does not link tangents into curves,
and the fingerprint image has no annotated bifurcations. The crossing test
therefore demonstrates retention of two directions, without establishing
branch recovery.

Additionally, the radius of 5 pixels limits how far support can bridge a gap. Endpoints
and image boundaries have fewer supporting neighbors and can lose confidence.

A single filter size cannot match both thin ridges and wide strands. The
template's negative center can respond to dark grooves and strand boundaries,
placing vectors away from the intended centerline. The implementation omits
the paper's lateral maxima selection, leaving several parallel rows of
vectors across strands. These effects are especially visible in hair and
spaghetti.

The support reference uses an ideal horizontal line. It does not adjust for
the number of grid pixels encountered at other orientations, local contrast,
or missing neighbors at image boundaries. Updating orientations independently
allows crossings to survive, but can also strengthen several nearby or
incorrect directions at once.

\section{AI tools and reproducibility}
Codex assisted with code organization, debugging, simplification, parameter
tuning, figure generation, report editing, LaTeX compilation, and layout checks.


\begin{thebibliography}{9}
\bibitem{parent1989}
P. Parent and S. W. Zucker.
\textit{Trace Inference, Curvature Consistency, and Curve Detection}.
IEEE Transactions on Pattern Analysis and Machine Intelligence,
11(8):823--839, 1989.
\bibitem{hummel1983}
R. A. Hummel and S. W. Zucker.
\textit{On the Foundations of Relaxation Labeling Processes}.
IEEE Transactions on Pattern Analysis and Machine Intelligence,
PAMI-5(3):267--287, 1983.
\end{thebibliography}
\end{document}
